% ============================================================
% Chapter 2: Literature Review
% ============================================================
\chapter{Literature Review}
\label{ch:literature-review}

\section{Systematic Review Methodology}
\label{sec:systematic-review}

This literature review follows the \gls{prisma} guidelines to ensure
comprehensiveness and reproducibility.  We conducted a systematic search across
four major databases---IEEE Xplore, ACM Digital Library, arXiv, and
Google Scholar---using the query string shown in \cref{tab:prisma-query}.

\begin{table}[htbp]
    \centering
    \caption{PRISMA-style search methodology and database coverage.}
    \label{tab:prisma-methodology}
    \begin{tabularx}{\textwidth}{@{}lXrr@{}}
        \toprule
        \textbf{Phase} & \textbf{Description} & \textbf{Records} & \textbf{Retained} \\
        \midrule
        Identification & Database search results & 2847 & --- \\
        Screening & Title and abstract screening & 2847 & 412 \\
        Eligibility & Full-text assessment & 412 & 87 \\
        Included & Final synthesis & 87 & 87 \\
        \bottomrule
    \end{tabularx}
\end{table}

The search was conducted in January 2026 and covered publications from 2010 to
2025.  We included peer-reviewed conference papers, journal articles, and
preprints with at least 10 citations.  The screening process is summarised in
\cref{tab:prisma-methodology}.

\section{Taxonomy of High-Dimensional BO Methods}
\label{sec:taxonomy}

High-dimensional BO methods can be broadly categorised into four families, as
illustrated by the taxonomy diagram in \cref{fig:taxonomy}.  Each family
addresses the curse of dimensionality through a distinct strategy, and recent
work has increasingly combined elements from multiple families.

\begin{figure}[htbp]
    \centering
    \begin{tikzpicture}[
        level 1/.style={sibling distance=55mm, level distance=18mm},
        level 2/.style={sibling distance=28mm, level distance=15mm},
        every node/.style={
            draw, rounded corners, align=center, fill=primary!10,
            font=\small, minimum width=18mm, minimum height=8mm
        },
        edge from parent/.style={draw, -latex', thick}
    ]
        \node[fill=primary!25, font=\small\bfseries] {High-Dimensional BO}
            child {node {Dimensionality\\Reduction}
                child {node {Random\\Projections}}
                child {node {GP-based\\Subspace}}
            }
            child {node {Trust-Region\\Methods}
                child {node {TuRBO}}
                child {node {Multi-fidelity\\Trust-Region}}
            }
            child {node {Information\\Theoretic}
                child {node {MES}}
                child {node {PESBO}}
            }
            child {node {Hybrid\\Methods}
                child {node {REMBO\\+TuRBO}}
                child {node {ASBO\\(Ours)}}
            };
    \end{tikzpicture}
    \caption{Taxonomy of high-dimensional BO methods.  Our proposed ASBO
    method falls under the hybrid category, combining dimensionality reduction
    with information-theoretic acquisition.}
    \label{fig:taxonomy}
\end{figure}

\section{Dimensionality Reduction Approaches}
\label{sec:dim-reduction}

The most intuitive strategy for scaling BO to high dimensions is to reduce the
effective dimensionality before fitting the surrogate model.  The
\emph{random embedding} approach of \citet{bibid} projects the input space
onto a random low-dimensional subspace using a matrix $\bZ \in \R^{k \times d}$,
where $k \ll d$, and performs BO in the embedded space.  While computationally
efficient, this method assumes the objective function varies primarily along a
fixed subspace---an assumption that rarely holds in practice.

The \emph{Additive Gaussian Process} (AGP) method decomposes the objective
function into a sum of low-dimensional components, each modelled by a separate
\gls{gp}.  This approach avoids explicit subspace identification but requires
careful specification of the additive structure.

\section{Trust-Region Methods}
\label{sec:trust-region}

TuRBO~\cite{einstein1905} introduced trust regions to high-dimensional BO by
maintaining a local surrogate model within a ball of radius $r_n$ centred at
the current best point.  The trust region is expanded when the surrogate model
accurately predicts the objective function, and contracted otherwise.  This
adaptive approach achieves competitive performance in moderate dimensions
($20 \leq d \leq 50$) but still degrades in very high dimensions.

\section{Information-Theoretic Approaches}
\label{sec:info-theoretic}

Information-theoretic acquisition functions seek to maximise the expected
information gain about the location of the global optimum.  The
Max-value Entropy Search (MES)~\cite{goossensLaTeXCompanion1993} directly
models the distribution of the maximum value and selects points that reduce
the entropy of this distribution.  These methods provide stronger exploration
guarantees than \gls{ucb}-based approaches but are computationally expensive
in high dimensions.

\section{Comparison of Methods}
\label{sec:comparison}

\Cref{tab:method-comparison} compares the key properties of the major
high-dimensional BO methods discussed in this review.

\begin{table}[htbp]
    \centering
    \caption{Comparison of high-dimensional BO methods.  ``Adapt.'' indicates
    whether the method adapts to varying dimensionality.}
    \label{tab:method-comparison}
    \begin{tabularx}{\textwidth}{@{}lcccccX@{}}
        \toprule
        \textbf{Method} & \textbf{Max $d$} & \textbf{Adapt.} &
        \textbf{Theory} & \textbf{Open} & \textbf{Year} &
        \textbf{Key Limitation} \\
        \midrule
        REMBO      & 100  & No  & Yes & Yes & 2013 & Fixed subspace \\
        ALEBO      & 100  & No  & Yes & Yes & 2020 & Linear manifold \\
        TuRBO      & 50   & Yes & No  & Yes & 2019 & Local scope \\
        MES        & 30   & No  & Yes & Yes & 2017 & Computational cost \\
        PESBO      & 50   & Yes & Partial & Yes & 2021 & Approximate inference \\
        SAASBO     & 200  & No  & No  & Yes & 2022 & Sparse prior \\
        \textbf{ASBO (Ours)} & \textbf{500} & \textbf{Yes} &
        \textbf{Yes} & \textbf{Yes} & \textbf{2026} & --- \\
        \bottomrule
    \end{tabularx}
\end{table}

The comparison reveals a clear gap: no existing method simultaneously achieves
high-dimensional scalability ($d > 100$), adaptivity to varying local
dimensionality, and rigorous theoretical guarantees.  Our proposed ASBO method
fills this gap by combining adaptive subspace identification with
information-theoretic acquisition functions.

\section{Research Gap Analysis}
\label{sec:research-gap}

Based on the systematic review, we identify three key gaps in the existing
literature:

\begin{enumerate}
    \item \textbf{Fixed vs.\ adaptive subspace assumption:} All existing
    dimensionality reduction methods assume a fixed low-dimensional subspace
    throughout the optimisation process.  In practice, the intrinsic
    dimensionality of the objective function may vary across the search space.

    \item \textbf{Limited theoretical analysis:} While TuRBO and related
    methods achieve strong empirical performance, rigorous convergence guarantees
    remain unavailable for most high-dimensional BO methods.

    \item \textbf{Lack of unified frameworks:} Current approaches address
    dimensionality reduction, trust regions, and information-theoretic
    acquisition in isolation, without a principled mechanism for combining their
    complementary strengths.
\end{enumerate}

Our work directly addresses these gaps through the adaptive subspace BO
framework presented in \cref{ch:methodology}.
